% \input{\pSections "ssec-repro-mcaninch"}

%%  %%  %%  %%  %%
\subsection{McAninch Challenge}
%
\begin{frame}\frametitle{McAninch Challenge}
	%
Can we tell if this experiment was done in \bl{Albuquerque}?
	%
\end{frame}
%
%
\begin{frame}\frametitle{McAninch Challenge}
	%
\begin{enumerate}
	\item Albuquerque average \bl{elevation}: \href{https://en.wikipedia.org/wiki/Albuquerque,_New_Mexico}{1619 m}
	\item \bl{Boiling point} of water at 1600 m: $94.4^{\circ}$C
\end{enumerate}
%
\center
\ \\
How can we use this apparatus to determine if the \\ water is boiling at $94.4^{\circ}$C?
	%
\end{frame}
%
%
\begin{frame}\frametitle{McAninch Challenge}
	%
\center
Without refining the apparatus, how do we \bl{refine the solution}?
	%
\end{frame}
%
%
\begin{frame}\frametitle{McAninch Challenge}
\center
Let's insist that $100^{\circ}$ and $94.4^{\circ}$C be \bl{3 standard deviations} apart.\\[20pt]
	%
\pause
	%
	Eqn. \eqref{eq:bevington interpolation error numeric}: $\sigma_{x=10} = \paren{100-94.4}/3 = 5.6 / 3 = 1.9$
\end{frame}
%
%
\begin{frame}\frametitle{Solve for $\sigma_{0}$ and $\sigma_{1}$}
	%
Let's call the position error $\bl{0.01}$ cm.
%
\begin{equation}
\begin{array}{ccccc ccccc c}
	\sigma_{x=10}^{2} &=& \sigma_{0}^{2} &+& x^{2} & \sigma_{1}^{2} &+& a_{1}^{2} & \sigma_{x}^{2}, \\
	\sigma_{x=10}^{2} &=& \rd{\sigma_{0}}^{2} &+& 10^{2} & \rd{\sigma_{1}}^{2} &+& 10^{2} & \bl{0.01}^{2} &=& 1.9^{2}. \\
\end{array}
\label{eq:bevington interpolation error numeric}
\end{equation}
	%
\end{frame}
%
%
\begin{frame}\frametitle{Device Characterization}
	%
\begin{figure}[t]
	\begin{center}
			%
		\begin{overpic}[ scale = 0.9 ]
			{\pLocalGraphics bev/extrap-error}
			\put(30, -4)	{$measurements, m$}
			\put(-6, 18)	{\rotatebox{90}{$error\ at\ x=10$}}
		\end{overpic}
			%
	\end{center}
\end{figure}
	%
\end{frame}
%
%
\begin{frame}\frametitle{Measurement Goal $1.9$}
\begin{figure}[t]
	\begin{center}
			%
		\begin{overpic}[ scale = 0.9 ]
			{\pLocalGraphics bev/extrap-error-marked}
			\put(30, -4)	{$measurements, m$}
			\put(-6, 18)	{\rotatebox{90}{$error\ at\ x=10$}}
		\end{overpic}
			%
			%
	\end{center}
\end{figure}
\end{frame}
%
%
\begin{frame}\frametitle{Result}
	To use our code to \bl{distinguish altitude}, we must make a \bl{200 measurements}.
\end{frame}

\endinput  %  ==  ==  ==  ==  ==  ==  ==  ==  ==
